% ============================================================================
% Abstract (TBD — was empty in Arce's docx)
% Nature Methods uses a "bold first paragraph" (editorial summary, ~150 words)
% rather than a structured abstract.
% ============================================================================

\textbf{%
% TODO: Write ~150-word bold editorial paragraph summarising the method,
% the simulation validation, and the SHS application finding. Linda + Arce.
Environmental mixture exposures accumulate across the life course, yet existing statistical methods for evaluating their joint health effects are largely limited to single time points or do not admit a causal interpretation under time-dependent confounding. Here we introduce g-BKMR, a Bayesian semi-parametric method that combines Bayesian kernel machine regression with the g-formula to estimate nonlinear, non-additive joint causal effects of correlated time-varying exposures while handling time-dependent confounding and enabling component-wise variable selection. Through simulation studies we show that g-BKMR recovers the true dose-response curve across a range of data-generating scenarios and outperforms existing alternatives. Applying the method to the Strong Heart Study, we quantify the joint causal effect of six urinary metals measured at three visits on fasting blood glucose, identifying zinc and selenium as the components with the strongest signal and providing a concrete template for applied causal mixture analysis in environmental epidemiology.%
}
